To ask whether a small set of interactions is sufficient to explain how
day length sets the timing of the two clock-neuron groups, we built a
minimal mathematical model in which the morning (M) and evening (E) cells
are each described by a single phase, the position of that cell group in
its daily transcriptional cycle. Three influences act on these phases.
First, light draws both oscillators toward a common reference phase that
depends on day length, so that with no other input M and E would peak
together. Second, a PDF-dependent drive, which acts in the light on days
longer than about 9~h, pushes the two apart by advancing M and delaying E.
Third, a weak mutual coupling pulls M and E back toward each other.
Because the PDF-dependent drive is equal and opposite in the two cell
groups, it leaves their average timing unchanged and serves only to open
the M--E gap; the width of the gap is set by the balance between this
separating drive and the two restoring influences, light and coupling.
With four fitted parameters, the model reproduced the measured
\textit{tim} peak times of M and E at all 11 photoperiods
(root-mean-square error, 0.23~h): coincident peaks on short days, a gap
that opens above 9~h of light and widens to about 4~h, and saturation of
both peaks on days longer than 15~h. The same parameters captured the
effects of removing PDF signalling at 12L:12D. Loss of PDF collapsed the
two peaks onto a single intermediate phase, whereas loss of the PDF
receptor in M cells alone or in E cells alone shifted only the affected
population and roughly halved the gap. The mutual coupling allowed the gap
to close progressively in constant darkness, from 2.9~h on the first day
to 1.2~h on the fourth, slightly short of the convergence to within about
1~h that we used as the experimental benchmark. The fit constrains the
ratio of the PDF-dependent drive to the light input more tightly than
either strength alone, as expected if the gap reports their balance.
Together, these results support competition between a PDF-dependent
separating drive and synchronizing influences as a parsimonious
explanation of the seasonal phase organization. The model takes day
length as an input, and so does not address how the clock neurons measure
it.
